\section{Equality and a common feasibility witness}
\label{sec:zero-deficit}

The integral in \eqref{eq:simplex-integral} measures how much of the
polar body is covered by feasible simplices. For a body attaining the
sharp volume product, we shall make the uncovered volume tend to zero
after adjoining one segment. Compactness then gives a feasible witness
at every interior point and in every direction. The resulting
representation will be the input to the equality classification.

Throughout this section $K=-K\subset\mathbb R^n$ is a convex body,
$n\ge1$, satisfying
\[
 |K|\,|K^\circ|=\frac{4^n}{n!}.
\]
Set $Q=K^\circ$, $d=n+1$, and
\begin{equation}\label{eq:segment-lift}
 B=K\oplus_1[-1,1]
   =\{(x,s):p_K(x)+|s|\le1\},
 \qquad C=B^\circ=Q\times[-1,1].
\end{equation}
The polar $C$ has the two horizontal faces $Q\times\{1\}$ and
$Q\times\{-1\}$. The representations constructed below will use
vectors from these faces.
In this lifted space, $X$ and $Y$ belong to $\mathbb R^d$, while elements of $C$
act on them by the Euclidean pairing. Recall that the lens width
$\lambda:[-1,1]\to[0,\infty)$ is continuous, even and concave, with
$\lambda(0)>0$ and $\lambda(t)\le\lambda(0)$.

\begin{proposition}[A representation minimizing the lens cost]
\label{prop:optimal-representation}
Let $K\subset\mathbb R^n$, $n\ge1$, be an origin-symmetric convex body
with $P(K)=4^n/n!$. Put $Q=K^\circ$, $d=n+1$, and define the lifted
bodies $B,C$ by \eqref{eq:segment-lift}.
For every $X\in\operatorname{int}B$ and every nonzero
$\xi=(q,a)\in\mathbb R^n\times\mathbb R$, put
\[
 T=p_C(\xi)=\max\{p_Q(q),|a|\}.
\]
There are $1\le r\le d$ vectors $c_i=(b_i,\sigma_i)$ with
$b_i\in Q$ and $\sigma_i\in\{-1,1\}$, nonnegative numbers
$\alpha_i$, and a vector $Y\in\mathbb R^d$, such that
\begin{equation}\label{eq:optimal-representation}
 \xi=\sum_{i=1}^r\alpha_i c_i,
 \qquad \sum_{i=1}^r\alpha_i=T,
\end{equation}
and
\begin{equation}\label{eq:common-witness}
 |h\cdot Y|\le\lambda(h\cdot X)\quad(h\in C),
 \qquad
 c_i\cdot Y=\lambda(c_i\cdot X)\quad(1\le i\le r).
\end{equation}
In particular, for every finite representation
$\xi=\sum_{j=1}^s\beta_jh_j$ with $\beta_j\ge0$ and $h_j\in C$,
\begin{equation}\label{eq:optimal-cost}
 \sum_{i=1}^r\alpha_i\lambda(c_i\cdot X)
 \le \sum_{j=1}^s\beta_j\lambda(h_j\cdot X).
\end{equation}
The choices may depend on $X$ and $\xi$; the same choice satisfies
\eqref{eq:optimal-cost} for all the competing representations.
\end{proposition}

The last coordinates $\sigma_i$ split the chosen rows into two groups.
In Section~\ref{sec:median}, we compare representation costs as $X$ approaches
the apex $(0,1)$ of $B$; the endpoint behavior of $\lambda$ makes the
two signs contribute differently. We now prove
Proposition~\ref{prop:optimal-representation} in three steps: approximate
$C$ by polytopes, locate feasible simplices near each prescribed point
and direction, and pass to a common witness to compare costs.

\subsection{Inner polytopes and vanishing missing volume}

The horizontal section of $B$ at height $s\in[-1,1]$ is
$(1-|s|)K$. Hence
\[
 |B|=|K|\int_{-1}^1(1-|s|)^n\,ds
     =\frac{2|K|}{n+1},
 \qquad |C|=2|Q|,
\]
and therefore
\begin{equation}\label{eq:lift-equality}
 |B|\,|C|=\frac4{n+1}|K|\,|Q|=\frac{4^d}{d!}.
\end{equation}

Choose full-dimensional symmetric polytopes $Q_j\subset Q$
converging to $Q$ in Hausdorff distance. For example, take the convex
hulls of increasingly fine finite symmetric nets in $Q$, including
a fixed spanning symmetric set. There are numbers $t_j\in(0,1)$
with $t_j\to1$ such that
\[
 t_jQ\subseteq Q_j\subseteq Q.
\]
Indeed, if $Q$ contains a Euclidean ball of radius $r_0>0$ and the
Hausdorff error is $\varepsilon_j$, then on Euclidean unit directions
\[
 h_{Q_j}\ge h_Q-\varepsilon_j
          \ge (1-\varepsilon_j/r_0)h_Q.
\]
The support-function criterion for containment gives the claim,
after discarding finitely many indices and choosing $t_j$ slightly
smaller if necessary.

Put $C_j=Q_j\times[-1,1]$ and $B_j=C_j^\circ$. The inclusions above
give
\begin{equation}\label{eq:lift-sandwich}
 t_jC\subseteq C_j\subseteq C,
 \qquad
 B\subseteq B_j\subseteq t_j^{-1}B.
\end{equation}
In particular, $|B_j|\,|C_j|\to|B|\,|C|$.
Use one row from each antipodal vertex pair of $C_j$ to describe $B_j$
as an intersection of strips. These rows span $\mathbb R^d$, and
no two are proportional: a ray from the interior point $0$ meets the
boundary of a convex body in only one point. Every signed row is a
vertex of $C_j$, and thus belongs to $Q_j\times\{-1,1\}$.

Let $\Sigma_{j,X}$ be the union of the feasible signed basis
simplices for this strip description, as in
\eqref{eq:simplex-integral}. In particular, for every
$X\in\operatorname{int}B_j$ it is a finite union of closed simplices
contained in $C_j$. Its definition is retained also at the exceptional
interior points where several constraints are simultaneously tight.
The integral estimate and \eqref{eq:lift-equality} imply
\begin{equation}\label{eq:vanishing-missing-volume}
 \begin{split}
 0\le D_j
 &:=\int_{B_j}\bigl(|C_j|-|\Sigma_{j,X}|\bigr)\,dX\\
 &\le |B_j|\,|C_j|-\frac{4^d}{d!}
 \longrightarrow0.
 \end{split}
\end{equation}

\subsection{From an integral deficit to each point and direction}

Fix $X\in\operatorname{int}B$ and $e\in\partial C$. We first
construct indices $j_k\to\infty$ and points
\begin{equation}\label{eq:nearby-feasible-pairs}
 X_k\in\operatorname{int}B_{j_k},\qquad
 e_k\in\Sigma_{j_k,X_k},\qquad
 (X_k,e_k)\longrightarrow(X,e).
\end{equation}
This is a consequence of \eqref{eq:vanishing-missing-volume} on
product neighborhoods; it does not require a pointwise limit of the
functions $X\mapsto|\Sigma_{j,X}|$.

For each $k$, choose a Euclidean ball $U_k$ of positive volume
contained in $\operatorname{int}B\cap B(X,1/k)$. There is also a ball
$V_k$ of positive volume with closure contained in
$\operatorname{int}C\cap B(e,1/k)$. To see this at the boundary point
$e$, first move a small distance from $e$ toward $0$; the resulting
point is interior because $0\in\operatorname{int}C$. Take a sufficiently
small ball about that point. By \eqref{eq:lift-sandwich},
$U_k\subset\operatorname{int}B_j$ for every $j$, and
$V_k\subset C_j$ for all sufficiently large $j$. The second assertion
follows also from $\max_{h\in\overline{V_k}}p_C(h)<1$ and $t_j\to1$.

Choose $j_k>j_{k-1}$ so large that these inclusions hold and
$D_{j_k}<|U_k|\,|V_k|$. If $\Sigma_{j_k,Z}$ missed $V_k$ for every
$Z\in U_k$, its missing volume would be at least $|V_k|$ throughout
$U_k$, contradicting this inequality. Thus choose
$X_k\in U_k$ and $e_k\in\Sigma_{j_k,X_k}\cap V_k$.
This proves \eqref{eq:nearby-feasible-pairs}.

Choose one feasible simplex containing $e_k$. Its $d$ signed basis
rows, written $c_{1k},\ldots,c_{dk}\in Q_{j_k}\times\{-1,1\}$,
give weights $\theta_{ik}\ge0$ satisfying
\begin{equation}\label{eq:approximate-representation}
 e_k=\sum_{i=1}^d\theta_{ik}c_{ik},
 \qquad \sum_{i=1}^d\theta_{ik}\le1.
\end{equation}
The coefficient left over is the weight of the simplex vertex $0$.
Feasibility supplies one vector $Y_k$ such that
\[
 c_{ik}\cdot Y_k=\lambda(c_{ik}\cdot X_k)\quad(1\le i\le d),
 \qquad
 |v\cdot Y_k|\le\lambda(v\cdot X_k)
 \quad\text{for every vertex $v$ of $C_{j_k}$}.
\]
Here a sign on a basis row has been absorbed using the evenness of
$\lambda$.

The last inequality extends to all $h\in C_{j_k}$. Indeed, write
$h=\sum_\ell s_\ell v_\ell$ as a convex combination of vertices.
The triangle inequality followed by concavity gives
\[
 |h\cdot Y_k|
 \le\sum_\ell s_\ell|v_\ell\cdot Y_k|
 \le\sum_\ell s_\ell\lambda(v_\ell\cdot X_k)
 \le\lambda(h\cdot X_k).
\]
Consequently $Y_k\in\lambda(0)B_{j_k}$. The polar sandwich
\eqref{eq:lift-sandwich} makes these witnesses uniformly bounded,
even if the selected bases approach singularity.

Pass to a subsequence on which $Y_k\to Y$, every $c_{ik}\to c_i$,
and every $\theta_{ik}\to\theta_i$. This uses a fixed finite product
of compact sets: all rows belong to $Q\times\{-1,1\}$ and all
weights belong to $[0,1]$. No independence of the limiting rows is
needed. We have
\[
 c_i\in Q\times\{-1,1\},\qquad
 e=\sum_{i=1}^d\theta_i c_i,\qquad
 \theta_i\ge0,\qquad \sum_i\theta_i\le1,
\]
and continuity of $\lambda$ gives
$c_i\cdot Y=\lambda(c_i\cdot X)$ for all $i$.
For arbitrary $h\in C$, use $h_k=t_{j_k}h\in C_{j_k}$ in the
feasibility inequality and pass to the limit. It follows that
\begin{equation}\label{eq:limiting-feasibility}
 |h\cdot Y|\le\lambda(h\cdot X)\qquad(h\in C).
\end{equation}
The subsequence has already been fixed; this argument proves
\eqref{eq:limiting-feasibility} for all $h$ with the same $Y$.

Since $e\in\partial C$, its gauge is one. Therefore
\[
 1=p_C(e)\le\sum_{i=1}^d\theta_i p_C(c_i)
   \le\sum_{i=1}^d\theta_i\le1.
\]
Thus $\sum_i\theta_i=1$: none of the limiting mass remains at the
simplex vertex $0$.

\subsection{Comparison with every positive representation}

To prove Proposition~\ref{prop:optimal-representation}, fix its
$X$ and nonzero $\xi$, and apply the preceding construction to
$e=\xi/T\in\partial C$. Set $\alpha_i=T\theta_i$ and discard any
zero weights if desired. This gives at most $d$ rows and proves
\eqref{eq:optimal-representation} and \eqref{eq:common-witness}.
For any finite competing representation
$\xi=\sum_j\beta_jh_j$, $\beta_j\ge0$, $h_j\in C$, the common
witness gives
\[
 \sum_i\alpha_i\lambda(c_i\cdot X)
 =\xi\cdot Y
 =\sum_j\beta_j h_j\cdot Y
 \le\sum_j\beta_j\lambda(h_j\cdot X).
\]
This proves \eqref{eq:optimal-cost} and completes the proof of
Proposition~\ref{prop:optimal-representation}. In particular,
the conclusion holds for each prescribed $X$ and $\xi$, including
directions for which every limiting representation is degenerate.
